\scanpage{56}
% archive foliation: 55
fibrés de $=$ espèce sur $S$, sous réserve que $T\to S$ soit homogène.

En particulier, si $G$ est un groupe fini ordinaire, un \deleted{\unsure{revêtem.}} revêtement principal de groupe $G$ de $Y/T$ est transformé en un \deleted{revêtement} fibré principal \unsure{lisse} sur $f_*(Y/T)$, de groupe $f_*(Y/T)\times_S f_*(\deleted{\unsure{G\times_S}}T\times G/T)$.
\margin{\unsure{Dans le résumé} on \unsure{supposera} $T/S$ fini, \unsure{bien que} \unsure{si on sait que la démonstration est générale}.}

Cela montre l'intérêt de déterminer, pour un ens. fini $I$ \deleted{\unsure{non vide}}, la variété \deleted{\unsure{un}} étale $f_*(T\times I/T)=\varphi(I)$\deleted{\unsure{étant}}.

\deleted{
On sait que $\varphi(T)(I)$ est un foncteur \deleted{fini} \uline{multiplicatif}, qui va des ens. finis aux variétés étales de $S$. Choisissons un pt géométrique $a$ sur $S$ supposé connexe, on obtient ainsi un foncteur $I\mapsto\varphi_a(T)(I)$ (Ens finis) dans les ens \inserted{(finis)} où $\pi_1(S,a)$ opère. En tant que foncteur à valeurs dans les ens finis, il est de la forme $\operatorname{Hom}_{\mathrm{Ens}\,f}(A_a(T),I)$, [petit raisonnement], et comme $\pi_1(S,a)$ opère dans ce foncteur, il opère dans $A_a(T)$, \deleted{\unsure{qui}} on a ainsi défini
\margin{et $=$ \unsure{exact}: produits}
\margin{\unsure{prendre pour} \unsure{principe}}
}
